\documentclass[a4paper, 11pt]{scrartcl}
\usepackage[left=3cm, right=3cm, top=3cm]{geometry}
\usepackage[utf8]{inputenc}
\usepackage[T1]{fontenc}
\usepackage{etoolbox}
\usepackage{microtype}
\usepackage{xcolor}
\usepackage[colorlinks=true,linkcolor=blue, allcolors=blue]{hyperref}%
\usepackage[normalem]{ulem}
% \pagenumbering{gobble}
\setkomafont{disposition}{\normalfont\bfseries}
\hypersetup{pageanchor=false}

\begin{document}

\title{\vspace{-2cm} \LARGE ALIUS Research Group}
\subtitle{\Large Internal Rules}
\author{}
\date{\large \today}
\maketitle

\section{Composition of ALIUS}
\subsection{Types of memberships}

ALIUS includes two types of members: \emph{research members} and \emph{non-research members}.

\begin{itemize}

\item \emph{Research members} are academic researchers whose work tackles (exclusively or non-exclusively) issues related to consciousness and the diversity of conscious states (see the inclusion criteria in \hyperref[new_memberships]{section 3.1}). The number of research members is unlimited in principle, but constrained in practice by membership renewals (see \hyperref[membership_renewal]{section 3.2}). 

\item \sout{The group can have a maximum of 3 \emph{non-research members} at any one time.} The role of \sout{these} \textcolor{red}{non-research} members is to help with outreach and communication (e.g., on \href{https://www.facebook.com/aliusresearch/}{social media}), with infrastructure (e.g., the \href{https://www.aliusresearch.org/}{website}) and with specific activities like the \href{https://www.aliusresearch.org/bulletin.html}{ALIUS Bulletin} or the ALIUS Podcast.


\end{itemize}

\noindent All ALIUS members will be involved in all decisions concerning the organisation and functioning of the group, through the following actions:

\begin{enumerate}
\item Electing coordinators every three years (see \hyperref[coordinators]{section 1.2}).
\item Proposing amendments to the present rules, and voting on such amendments twice a year (see \hyperref[modif]{section 2}).
\item Voting on new memberships twice a year (see \hyperref[new_memberships]{section 3.1})

\end{enumerate}

\subsection{Coordinators}
\label{coordinators}

Three members of ALIUS will be elected to become coordinators of the group. Coordinators are in charge of administrative tasks within the group (such as collecting motions and organizing votes), as well as helping with communication on social media. They are responsible for enforcing the policies outlined in this document, including respecting the inclusion criteria for new members (\hyperref[new_memberships]{section 3.1}) and the criteria for the renewal of existing memberships (\hyperref[membership_renewal]{section 3.2}).

More specifically, the roles of coordinators include the following:

\begin{enumerate}

\item Organizing a vote on the invitation of new members to join the group every six months (if there has been any suggestion of people to invite during this timeframe).
\item Ensuring that members whose membership is renewed (every three years, see below) write a short `note of intent' presenting their research and activity within the group.
\item Promoting the group's activities and sharing relevant new publications of research members on the group's social media accounts. \sout{(Coordinators also have to update the list of publications on the website at least once every three years, but preferably more often.)}

\end{enumerate}

\noindent Coordinators are elected every three years by all members through an online vote. The preferred electoral system is the \href{https://en.wikipedia.org/wiki/Majority_judgment}{majority judgement}, which presents a number of advantages over \href{https://en.wikipedia.org/wiki/Condorcet_method}{Condorcet methods}. The vote will be organized on the online platform \href{https://jugementmajoritaire.net/}{jugementmajoritaire.net} (available in English) or equivalent.

\section{Modification of the rules}
\label{modif}

Any member can write at any time to the coordinators to propose a new motion to be discussed ahead of the following biannual vote. Such motions include any amendment of the present rules to modify or extend them. 

The coordinators will list all submitted motions up to one month prior the biannual vote by email. Every member can join the discussion by email to present arguments for or against each submitted motion. The anonymous vote will then be organized through an online platform. 

A motion will be adopted if the two following conditions are met at the end of the vote:

\begin{enumerate}
\item At least 40\% of members (rounded to the nearest integer) voted.
\item At least 70\% of votes (rounded to the nearest integer) are in favor of the motion.
\end{enumerate} 

\section{New memberships and membership renewals}

\subsection{New memberships}
\label{new_memberships}

Any existing member can write to the coordinators to propose to invite a new \emph{research member} or \emph{non-research member} in the group. They must motivate this proposition by providing a short presentation of the candidate's work. The coordinators must ensure that the proposition fulfills the following inclusion criteria: \\

For new \emph{research members}:

\begin{enumerate}
\item They must be affiliated to an academic institution. 
\item Their research must inform or be informed by the empirical study of consciousness and conscious states. 

\end{enumerate}


For new \emph{non-research members}:   

\begin{enumerate}

\item They must have some technical skill or communication skill that could be useful to develop the activities and online presence of ALIUS.

\end{enumerate}

\noindent The invitation of new members will be submitted to a vote every six months, provided that any proposition has been put forward by existing members. Research members and non-research members can take part in this vote.

A new member will be invited to join ALIUS after the vote if the two following conditions are met:

\begin{enumerate}
\item At least 25\% of members (rounded to the nearest integer) voted.
\item At least 80\% of votes (rounded to the nearest integer) are in favour of inviting the researcher.
\end{enumerate}

\noindent If the outcome of a vote is positive, coordinators will proceed to invite the chosen individual to become a member of the group.

ALIUS is committed to improving gender diversity within the group. To this end, the group must count at least a \emph{third} (33\% rounded to the nearest integer) of women among its members at any time.

\subsection{Renewal of existing memberships}
\label{membership_renewal}

The membership of existing research members will be evaluated and renewed every three years. The renewal of researchers' membership is conditional on the satisfaction of the following criteria within this three-year interval:

\begin{enumerate}

\item Having published at least one article or book in English on the topic of consciousness in a peer-reviewed journal of international reputation as a first author or corresponding author (for an article), or with a publisher of international reputation as a single author (for a book). Publications that have been written in another language but translated in English also qualify.

\sout{Having actively participated to at least \emph{two} of the following activities:}

\sout{Organizing a panel/workshop/conference that includes at least two ALIUS research members among the speakers}
\sout{Speaking at a panel/workshop/conference organized by an ALIUS researcher and including two ALIUS research members among the speakers}
\sout{Co-editing or contributing to the \href{https://www.aliusresearch.org/bulletin.html}{ALIUS Bulletin}}
\sout{Writing a post on one's research for the \href{https://www.aliusresearch.org/blog.html}{blog} on the ALIUS website (research members can write to the coordinators with a proposal for a short blog post at all times)}
 \sout{Contributing to the ALIUS Podcast (research members can write to the coordinators with a proposal for a podcast episode at all times)}
\sout{Helping with the administration of the group (e.g. by being a coordinator)}

\item \textcolor{red}{Having actively participated to at least once to the ALIUS Bulletin as a contributor or co-editor}

\sout{One month before the renewal date, writing a short `note of intent' (750 to 1500 words) presenting one's current research and explaining how one's work and collaboration with other ALIUS members (if any) is contributing to the group's mission of advancing research on consciousness. These `notes of intent' will be gradually published on the blog throughout the following year as a way to present the work of each research member.}
\item\textcolor{red}{At the renewal date, sharing how they intent to contribute to the ALIUS Bulletin within the 3 years to come.}

\end{enumerate}

If a research member fails to satisfy these three conditions at the time of their assessment after three years, their membership will not be renewed.

\begin{itemize}

\item For existing research members, the first membership renewal date will be three years from the date of their agreement with this document. 
\item For future research members, the first membership renewal date will be three years after the date of their inclusion within the group. 
\item Subsequent membership renewal dates for all research member will be three years after their last.
\item \textcolor{red}{If research activities are interrupted (e.g., maternity leave), the period of 3 years will be extended by the necessary amount of time.}
\end{itemize}

\end{document}


